\scanpage{33}
\begin{center}
Cours DEA 1977/78 de A. Grothendieck

\uline{Géométrie arithmétique des polygônes réguliers}
\end{center}

{\reversemarginpar\margin{\handwritten{!}}\normalmarginpar}%
% A short hand-drawn line under the typed words ``en commun''.
Ce ``cours'' se fera dans l'esprit d'une investigation \handwritten{\uline}{en commun} de multiples facettes de la notion de polygône régulier au-dessus d'un corps de base, puis d'un anneau de base quelconque. Ce sera une occasion stimulante pour nous familiariser avec les notions de base de géométrie vectorielle, affine, projective - avec ou sans métrique - en dimensions un, deux et trois, et avec les phénomènes particuliers à ces dimensions, ainsi qu'avec les aspects parfois imprévus que prennent des notions familières (comme celle de réflexion\deleted{\unsure{s}}, ou de centre d'un polygône régulier) sur les corps de caractéristique positive, et notamment sur les corps finis. Le besoin d'une géométrie sur des anneaux de base généraux (tel l'anneau ``sur lequel s'écrit'' le $\mathrm{n}$-polygône régulier ``le plus général'') se fera sentir en cours de route, et nous amènera sans doute à affleurer quelques points de fondements de géométrie algébrique. Si le temps le permet et si d'autres investigations ne nous retiennent , nous tenterons, dans la volée, de construire l'icosaèdre régulier ``général'' sur un anneau de base convenable, fini sur l'anneau ${\handwritten{\mathbb}{Z}}$ des entiers, ayant des corps résiduels de toutes caractéristiques. (La caractéristique deux \hspace{.5em}donnera sans doute du fil à retordre !).

A partir d'un point de départ intuitif et familier, ce ``cours'' voudrait être une ouverture à un univers fascinant, pleins de mystères qui ne demandent qu'à être élucidés. Il ne peut donner des fruits qu'avec une participation active de l'étudiant et un travail personnel important, notamment pour développer a\deleted{\unsure{o|e}}vec le détail nécessaire les nombreux points \handwritten{(}de fondement\inserted{\handwritten{s}} notamment) qui n'auront été qu'esquissés dans le ``cours''. Il lui est très fortement conseillé de suivre en même temps le cours C4 sur l'icosaèdre, qui se fera dans un esprit analogue. Il y aura sans doute de nombreuses interférences entre ces deux cours de réflexion.
% archive sheet number: 32
